\subsection{Perturbation of Riesz endomorphisms}

Let E be a Banach space. Recall (cf. III, p. 5, prop. 3) that the set \(\mathcal{L}^c (\mathrm{E})\) of compact endomorphisms of E forms a closed two-sided ideal of the Banach algebra \(\mathcal{L}(\mathrm{E})\) . The quotient Banach algebra is called the Calkin algebra of E; it is denoted \(\mathcal{Calk}(\mathrm{E})\) . We denote by \(\pi\) the canonical homomorphism of \(\mathcal{L}(\mathrm{E})\) onto \(\mathcal{Calk}(\mathrm{E})\) . By cor. 1 of III, p. 74, an endomorphism \(u\) of E is a Fredholm endomorphism if and only if \(\pi (u)\) is invertible in the algebra \(\mathcal{Calk}(\mathrm{E})\) .

PROPOSITION 2. --- Let \(u \in \mathcal{L}(\mathrm{E})\) such that \(\|1_{\mathrm{E}} - \pi(u)\| < 1\) in the algebra \(\mathcal{Calk}(\mathrm{E})\) . Then \(u\) is a Riesz endomorphism of E.

Let \(r \geqslant 1\) be an integer such that \(\| 1_{\mathrm{E}} - \pi(u)\|^{r} < \frac{1}{2}\) . Let \(\mathrm{P} \in \mathbf{C}[\mathrm{X}]\) be the polynomial \(\frac{1 - (1 - \mathrm{X})^r}{\mathrm{X}}\) . Set \(v = 1_{\mathrm{E}} - (1_{\mathrm{E}} - u)^r\) . We therefore have \(v = u\mathrm{P}(u)\) and \(\| 1_{\mathrm{E}} - \pi(v)\| < \frac{1}{2}\) . Since the endomorphisms \(u\) and \(\mathrm{P}(u)\) of E commute, it suffices to prove that \(v\) is a Riesz endomorphism of E (III, p. 49, prop. 9).

Since \(\|1_{\mathrm{E}}-\pi(v)\|<1/2\) , there exists, by definition of the quotient norm in the space \(\mathcal{Calk}(\mathrm{E})\) , a compact endomorphism h of E and an endomorphism w of E such that \(v=1_{E}+h+w\) and \(\|w\|<\frac{1}{2}\) . By corollary 1 of I, p. 22, the element \(1_{E}+w\) is an automorphism of E. Since h is compact, \(v=(1_{\mathrm{E}}+w)+h\) is a Fredholm endomorphism of E of index 0 (cor. 2 of III, p. 74). For every integer \(n\geqslant0\) , denote by \(N_{n}\) the kernel of \(v^{n}\) . To prove that v is a Riesz endomorphism of E, it suffices to show that there exists an integer \(n \geqslant 0\) such that \(N_{n} = N_{n+1}\) (III, p. 46, def. 2 and III, p. 46, remark).

Let us argue by contradiction by assuming the sequence \((\mathrm{N}_{n})\) strictly increasing. For every \(n \in N\) , denote by \(p_{n}\) the canonical map from E onto the normed space \(E/N_{n}\) . Let c be a real number such that \(2\|w\| < c < 1\) . Let \(n \in N\) . Since \(N_{n+1}\) is different from \(N_{n}\) , there exists an element \(x_{n} \in N_{n+1}\) such that \(\|p_{n}(x_{n})\| = c\) and such that \(\|x_{n}\| < 1\) (indeed, there exists \(y \in N_{n+1}/N_{n}\) of norm c, so for every \(\varepsilon > 0\) , there exist \(x_{n} \in N_{n+1}\) such that \(p_{n}(x_{n}) = y\) and \(\|x_{n}\| \leqslant c + \varepsilon\) ).

Let \(m\) , \(n\) be two integers such that \(m > n \geqslant 0\) . We have

\[
\begin{array}{r l} & {\| h (x _ {m}) - h (x _ {n}) \| = \| v (x _ {m} - x _ {n}) - (1 _ {\mathrm{E}} + w) (x _ {m} - x _ {n}) \|} \\ & {\qquad \geqslant \| (v - 1 _ {\mathrm{E}}) (x _ {m} - x _ {n}) \| - \| w (x _ {m} - x _ {n}) \|} \\ & {\qquad \geqslant \| (v - 1 _ {\mathrm{E}}) (x _ {m} - x _ {n}) \| - 2 \| w \|} \end{array}\tag{1}
\]

since \(\| x_m\| \leqslant 1\) and \(\| x_n\| \leqslant 1\) . Moreover, note that

\[
(v - 1 _ {\mathrm{E}}) (x _ {m} - x _ {n}) = v (x _ {m} - x _ {n}) + x _ {n} - x _ {m}.
\]

But since n \textless{} m and \(v(\mathrm{N}_{m+1}) \subset \mathrm{N}_{m}\) , we have \(v(x_{m} - x_{n}) + x_{n} \in \mathrm{N}_{m}\) , whence

\[
\left\| v (x _ {m} - x _ {n}) + x _ {n} - x _ {m} \right\| \geqslant \left\| p _ {m} (x _ {m}) \right\| = c.
\]

Inequality (1) therefore yields the lower bound

\[
\left\| h \left(x _ {m}\right) - h \left(x _ {n}\right) \right\| \geqslant c - 2 \| w \| > 0.\tag{2}
\]

The sequence \((h(x_{m}))_{m\in\mathbf{N}}\) therefore has no cluster point in E, which contradicts the fact that the image under h of the unit ball of E is relatively compact, since h is compact (remark 1 of III, p. 2).

COROLLARY 1. --- Let E be a separated locally convex space and let \(u \in \mathcal{L}(\mathrm{E})\) . For u to be a Riesz endomorphism, it is necessary and sufficient that there exist an element v of \(\mathcal{L}(\mathrm{E})\) which commutes with u and such that \(1_{\mathrm{E}} - u \circ v\) is compact.

Since every continuous linear map of finite rank is compact, the condition is necessary (III, p. 47, prop. 6 d)). Let us prove that it is sufficient. Let v be an element of \(\mathcal{L}(\mathrm{E})\) that commutes with u and such that the endomorphism \(h = 1_{E} - u \circ v\) of E is compact. There exist a Banach space G, a continuous linear map \(p: E \to G\) and a compact linear map \(q: G \to E\) such that \(h = q \circ p\) (III, p. 5, prop. 5). The endomorphism \(p \circ q\) of G is compact. By prop. 2, the linear map \(1_{G} - p \circ q\) is a Riesz endomorphism of G. But then \(u \circ v = 1_{E} - q \circ p\) is a Riesz endomorphism of E (III, p. 49, prop. 10, g)). Since u and v commute, u is a Riesz endomorphism of E (III, p. 49, prop. 9).

COROLLARY 2. --- Let E be a separated locally convex space, u a Riesz endomorphism of E, and \(h \in \mathcal{L}(\mathrm{E})\) a compact linear map. If u and h commute, then \(u + h\) is a Riesz endomorphism of E.

Let \(v\) be the canonical quasi-inverse of \(u\) . It commutes with \(u\) and with \(h\) (III, p. 47, prop. 6), hence with \(u + h\) . The endomorphism \(1_{\mathrm{E}} - (u + h) \circ v\) of E is the sum of the continuous linear map of finite rank \(1_{\mathrm{E}} - u \circ v\) and the compact linear map \(-h \circ v\) , hence is compact. It follows that \(u + h\) is a Riesz endomorphism of E (corollary 1).

PROPOSITION 3. --- Let E be a Banach space. Let u be an endomorphism of E, and let A be its commutant in \(\mathcal{L}(\mathrm{E})\) . The closure B of \(\pi(\mathrm{A})\) in the algebra \(\mathcal{Calk}(\mathrm{E})\) is a Banach algebra. For u to be a Riesz endomorphism of E, it is necessary and sufficient that \(\pi(u)\) be invertible in B.

Since \(\pi(A)\) is a unital subalgebra of \(\mathcal{Calk}(E)\) , its closure B is a Banach algebra. If \(u\) is a Riesz endomorphism of E, the element \(\pi(u)\) has an inverse in \(\pi(A)\) (cor. 1), hence in B. Conversely, suppose that \(\pi(u)\) has an inverse \(s\) in B. By definition of B, there exists an element \(v\) of A such that

\[
\left\| s - \pi (v) \right\| <   \frac {1}{\left\| \pi (u) \right\|},
\]

whence \(\|1_{E}-\pi(u\circ v)\|<1\) . By prop. 2, the linear map \(u\circ v\) is a Riesz endomorphism of E. The same is true of u since u and v commute (III, p. 49, prop. 9).

